\chapter{选题依据}
\section{实践背景：风险审核面对的是不断变化的信息集}
上市公司的短期风险筛查通常从一个很具体的问题开始：在当天可获取的资料中，哪些公司值得优先复核？审核人员可以看到近期收益、经营预告、诉讼进展和持股名单，却难以每天逐项阅读所有公司的材料。市场行情不断更新，公告出现的时间不规则，持股名单又带有明显的报告期和披露滞后。把这些资料整理为公司日风险排序，首先需要解决的并不是网络层数，而是资料在历史时点上的可用性。

业务日期和信息日期的区别，在低频披露中尤其重要。季度末描述的是持股状态对应的期间，公开日描述的是市场参与者何时能够获知该状态，数据库更新时间则可能只是更正或补录发生的时间。三者混用会产生两类问题。一类是把晚公布的内容提前放入回测；另一类是资料虽然确实已公开，却在很长时间内被当作同样新鲜的信息使用。前者影响实验有效性，后者影响模型对辅助信息的信任程度。图\ref{fig:async}把两类问题放在同一时间框架下。

\reportfigure{收益、公告与持股关系的异步更新}{fig:async}{\FigAsync}{研究流程示意。时间点仅表达先后次序，不是实际公司个案，也不表示经统计得到的更新频率。}

仅靠日期过滤不能解决全部问题。某份名单在预测日前已公开，仍可能因公司经营变化、机构调仓或名单更正而逐渐失去代表性。文本抽取也存在类似情形：模型可以输出一个格式正确的答案，答案却未必准确概括经营变化。若融合模型允许这类资料大幅改变主干分数，新增信息可能不但没有帮助，还会增加排序波动。因此，资料的存在、可追溯程度和陈旧程度需要进入模型设计，而不是只保存在数据清洗日志里。

本课题面向上述审核场景，研究A股公司未来五个市场交易日的条件左尾事件。主干使用持续更新的收益和三因子序列；业绩预告语义、共同机构持股关系作为异步辅助资料；诉讼信息形成可直接追溯的结构控制。研究输出是离线风险排序方法和原型，供有限审核容量下的风险筛查研究使用。模型不判断某条公告是否因果性地引起下跌，也不把共同持股解释为已经识别的风险传染机制。

\section{研究对象与任务边界}
“尾部”需要在任务中给出可以重复计算的定义。普通涨跌分类把零收益作为分界，而年度崩盘研究常使用较长时期的偏度或崩盘代理变量。本课题选择未来五日累计因子残差相对公司自身历史分布的左尾。这样能够将近期价格状态和公司特有异常波动联系起来，也允许不同公司使用适应自身历史分布的阈值。它不是监管定义的风险等级，更不是破产概率。

预测时点定为交易日收盘。收益、因子及公告关系的查询都以该时点为界，未来收益只用于离线定义结果标签。没有可靠时分秒的披露资料，从公开后的下一市场交易日起进入特征。采用这一约定会损失部分当天已公开的信息，但可以避免把盘后材料当作盘中已知事实。对于目前以日期为粒度的源数据，这是一项明确、保守且可以程序验证的取舍。

图\ref{fig:availability}还区分了更正内容。更正不是把原始记录的一个数值“修好”以后再覆盖所有历史日期，而是新的可见版本。训练与测试的历史查询必须返回当时能够使用的内容。若源表无法提供旧版本，原公告比对和数据隔离就成为必要环节。资料认证的结果会影响辅助模态覆盖，但不应改变同一个公司日标签是否参加模型比较。

\reportfigure{报告期、公开日、生效日和更正版本的关系}{fig:availability}{\FigAvailability}{按日期级披露拟定的可用性规则。更正保留独立版本，不据报告期回填。}

\section{相关研究为本课题提供了什么}
\subsection{崩盘与公司特有收益}
Chen、Hong和Stein从条件偏度讨论股票崩盘预测，强调过去市场状态与极端下跌分布之间的联系\cite{r1}。Jin和Myers关注信息不透明、收益共同波动及风险积累\cite{r2}，Hutton等将财务报告不透明性与崩盘风险联系起来\cite{r3}。这些研究使公司特有信息成为尾部研究的重要线索，但其年度变量、经济解释和样本结构与本课题的日频预测不同。

Fama和French三因子模型提供了一个熟悉且可计算的共同风险参照\cite{r4}。这里使用三因子并非认为它已经解释全部收益变化，而是为标签建立一致口径：预测起点估计系数，冻结系数后观察未来残差，再与过去已结束窗口比较。这一操作定义有利于复核，但残差仍会包含遗漏因子、估计误差和价格噪声。实验结论应始终对应这一标签，不把它扩写为“识别所有形式的金融风险”。

\subsection{文本从情绪走向事件与证据}
Tetlock及其后续研究说明新闻语言可能包含数值变量之外的信息\cite{r5,r6}。Loughran和McDonald指出财务语境会改变词语的含义\cite{r7}。这些线索对业绩预告尤其直接：盈利但同比下降、仍亏损但亏损收窄，分别涉及绝对状态和变化方向，不能只根据“盈利”或“亏损”两个词判断。

EFSA将中文财经新闻分析细化到公司、行业、事件和情感的组合\cite{r24}。本课题借鉴事件级分析思路，但把正式预告的抽取目标限定为方向、明确原因和原文证据。选择有限字段，一方面使人工规范能够执行，另一方面便于追查模型为什么产生某个事件表示。预训练模型、LoRA与QLoRA提供的是表示和适配工具\cite{r14,r15,r16,r21,r22,r23}；是否适合本地公告，还需要独立金标评价。

\subsection{关系网络与融合模型的直接近邻}
Feng等研究股票关系下的时序排序\cite{r8}，Cheng等研究多模态图神经网络用于金融序列\cite{r9}，Hou和Liu直接讨论机构交叉持股与崩盘风险\cite{r10}。关系与风险的联系因此不是空白。GraphSAGE、图卷积和图注意力提供了多种邻居聚合方式\cite{r17,r18,r19}，但这些算子并不自动解决持股名单来源和金融主体身份的问题。

Jing等把改造Graph WaveNet用于中国股票崩盘风险\cite{r11,r20}；MSGCA同时利用指标、动态文本和图关系，并设置门控与主干回接\cite{r12}。两者要求本课题分别面对任务方向和融合架构上的近邻，而不能仅与一个弱线性分类器比较。表\ref{tab:neighbors}将这种关系具体化。不同论文的频率、标签和信息集并不一致，原论文分数不能直接排入本地结果表。

\reporttable{主要近邻与本任务需要重新验证的部分}{tab:neighbors}{
\begin{tabularx}{\linewidth}{p{2.65cm}YY}\toprule
研究线索&已有内容&本课题中的比较问题\\\midrule
经典崩盘研究&条件偏度、信息透明度、公司治理&年度解释性变量能否转化为只用历史的日频特征\\
文本与EFSA&财经语言和事件级语义&正式预告三字段是否抽取准确、是否有下游增量\\
共同持股&机构交叉持股与崩盘风险联系&公开时点图是否超过收益与公告信息\\
Graph WaveNet&图时序建模及中国股票崩盘应用&同一五日标签、同一时间折的任务适配表现\\
MSGCA&指标、文本、图的门控融合&同编码器下年龄幅度约束是否仍有必要\\\bottomrule
\end{tabularx}}

\section{本课题的研究切入点}
相近文献已经拥有三路输入、门控、关系建模和收益主干。本课题的切入点不是再增加一个组件，而是把辅助信息对主干分数的作用限制为可用性、质量、年龄和有界内容项的乘积。这种设计可以先给出修正幅度上界，再在相同信息条件下检验其代价和收益。

关键问题是约束是否过强。稳定的持股关系可能在较长时间里仍有效，一份旧公告也可能解释持续经营变化。年龄越大不等于信息必然无用，所以衰减率需要由训练学习，同时设置年龄仅作普通特征、去全部年龄条件和仅去边级衰减的控制。它们可以区分“增加年龄信息”与“限制修正幅度”，也能检验图内相对权重和图外绝对陈旧度分别起什么作用。

专业学位研究还需要把方法落实到可以检查的工程对象。公司编码、公开版本、正文缓存和日批次索引都应有清楚的接口与失效条件。其价值在于让风险排序能追溯到具体资料与模型版本，而非只保留一个分数。本文后续据此展开数据、模型、实验和实现设计；独立文献综述另行提交，不把本章当作对完整综述的重复粘贴。
